% !TEX program = xelatex
% 공통수학2 · Ⅲ. 함수와 그래프 · 2. 유리함수와 무리함수 · 01 유리함수 · 1/2차시
\documentclass[11pt,a4paper]{article}
\usepackage{kotex}
\usepackage{iftex}
\ifPDFTeX\else
  \setmainhangulfont{Noto Serif CJK KR}
  \setsanshangulfont{Noto Sans CJK KR}
\fi
\usepackage[margin=2.1cm]{geometry}
\usepackage{parskip}
\usepackage{enumitem}
\usepackage{array}
\usepackage{tabularx}
\usepackage{booktabs}
\usepackage{xcolor}
\usepackage{amssymb}
\usepackage{tikz}
\usepackage[most]{tcolorbox}
\usepackage{fancyhdr}
\usepackage{titlesec}

\definecolor{goldc}{HTML}{B07C22}
\definecolor{goldstrong}{HTML}{8A5F16}
\definecolor{indigoc}{HTML}{2B3B57}
\definecolor{indigosoft}{HTML}{6E82A6}
\definecolor{linec}{HTML}{D6DACB}
\definecolor{surface2}{HTML}{E4E8DC}
\definecolor{notebg}{HTML}{FBF3E3}
\definecolor{inksoft}{HTML}{52604F}

\titleformat{\section}{\normalfont\sffamily\bfseries\large\color{indigoc}}{\thesection}{0.6em}{}
\titlespacing{\section}{0pt}{16pt}{8pt}
\newtcolorbox{ideabox}{enhanced, breakable, colback=white, colframe=goldc, boxrule=0.5pt, leftrule=3pt, arc=2pt, left=10pt, right=10pt, top=8pt, bottom=8pt}
\newtcolorbox{calloutbox}{enhanced, breakable, colback=white, colframe=indigosoft, boxrule=0.5pt, leftrule=3pt, arc=2pt, left=10pt, right=10pt, top=7pt, bottom=7pt}
\newtcolorbox{notebox}{enhanced, breakable, colback=notebg, colframe=goldc, boxrule=0.5pt, leftrule=3pt, arc=2pt, left=10pt, right=10pt, top=7pt, bottom=7pt}
\newtcolorbox{answerbox}{enhanced, breakable, colback=white, colframe=linec, boxrule=0.5pt, arc=2pt, left=10pt, right=10pt, top=7pt, bottom=7pt}
\newcommand{\lessonbanner}[3]{%
  \begin{tcolorbox}[enhanced, colback=indigoc, colframe=indigoc, boxrule=0pt, arc=0pt,
    left=14pt, right=14pt, top=14pt, bottom=10pt, borderline south={2.4pt}{0pt}{goldc}, fontupper=\color{white}]
    {\sffamily\footnotesize\color{white!85!black} #1}\\[4pt]{\sffamily\bfseries\Large #2}\\[3pt]{\sffamily\small\color{white!88!black} #3}
  \end{tcolorbox}}
\pagestyle{fancy}\fancyhf{}\renewcommand{\headrulewidth}{0pt}
\fancyfoot[L]{\footnotesize\color{inksoft} 공통수학2 · 2026학년도 2학기 · 지평선고등학교}
\fancyfoot[R]{\footnotesize\color{inksoft} 교사용 지도안 -- 배포 금지}

\begin{document}
\lessonbanner{공통수학2 · Ⅲ. 함수와 그래프 · 2. 유리함수와 무리함수}{01 유리함수 --- 1/2차시}{유리함수의 뜻과 정의역 · 교사용 지도안 (2026학년도 2학기)}
\vspace{10pt}

\noindent\begin{tabularx}{\linewidth}{@{}>{\bfseries\color{indigoc}}p{2.6cm}X@{}}
\toprule
대상 & 고1 · 2개 반 · 반당 13명 \\
차시 & 50분 (도입 10 · 전개 30 · 정리 10) \\
성취기준 & 유리함수의 뜻을 알고, 정의역을 구할 수 있다. \\
선행 학습 & 38차시까지 --- `1. 함수' 전체(함수, 합성함수, 역함수) \\
\bottomrule
\end{tabularx}

\section{학습 목표 \& Big Idea}
\begin{ideabox}
\textbf{\color{goldstrong}영속적 이해 (Big Idea)}\\
유리함수는 분모가 $0$이 되는 지점에서 나눗셈 자체가 의미를 잃기 때문에, 그 지점을 제외한 곳에서만 정의된다. 이 `제외된 지점'은 예외가 아니라 유리함수의 본질적인 성질이다.
\vspace{4pt}
\small\color{inksoft} 핵심개념: \textbf{분모의 제한} \quad\textbullet\quad 핵심개념: \textbf{식의 변형} \quad\textbullet\quad 일반화: \textbf{나눌 수 없는 지점이 함수의 모양을 결정한다}
\end{ideabox}
\begin{calloutbox}
\textbf{차시 학습 목표}
\begin{itemize}[leftmargin=1.4em, itemsep=2pt, topsep=4pt]
  \item 유리식과 유리함수의 뜻을 알고, 다항함수도 유리함수의 특수한 경우임을 설명할 수 있다.
  \item 유리함수의 정의역을 구하고, $y=\dfrac{ax+b}{cx+d}$ 꼴을 $y=\dfrac{k}{x-p}+q$ 꼴로 변형할 수 있다.
\end{itemize}
\end{calloutbox}

\section{질문의 연결}
\begin{tabularx}{\linewidth}{@{}>{\bfseries\color{indigoc}\small}p{2.4cm}X@{}}
사실적 질문 & 유리함수 $y=\dfrac{ax+b}{cx+d}$의 정의역은 어떻게 구할까? \\[6pt]
개념적 질문 & 다항함수 $y=2x+1$도 유리함수라고 할 수 있을까? 왜 그럴까? \\[6pt]
논쟁적 질문 & 피자 한 판을 $n$명이 똑같이 나눠 먹을 때, $n=0$이면 `한 사람 몫'은 어떤 의미를 가질까? 수학이 `정의되지 않는다'고 말하는 것과 현실에서 `말이 안 된다'고 하는 것은 같은 것일까? \\
\end{tabularx}

\section{수업 흐름}
\begin{tabularx}{\linewidth}{@{}>{\centering\arraybackslash}p{1.6cm}X>{\raggedright\arraybackslash\small}p{3.1cm}@{}}
\toprule
\textbf{단계} & \textbf{활동 \& 발문} & \textbf{ATL / VTR} \\
\midrule
\textcolor{goldstrong}{\textbf{도입}}\newline\footnotesize 10분 &
\textbf{마음공부 (2분)} --- 새 소단원 시작, 유무념 대조.\newline
\textbf{생각 열기 (8분)} --- 피자 한 판을 $n$명이 나눌 때 한 사람 몫 $\frac{1}{n}$을 표로 만들며 $n$이 $0$에 가까워질 때 값이 어떻게 변하는지, $n=0$일 때 왜 정의할 수 없는지 토의. &
ATL: 사고(극한적 직관)\newline VTR: See-Think-Wonder \\
\addlinespace
\textcolor{goldstrong}{\textbf{전개}}\newline\footnotesize 30분 &
\textbf{유리식과 유리함수 (10분)} --- 다항식/다항식 꼴이 유리식, 유리식으로 나타낸 함수가 유리함수임을 정의. 다항함수(분모가 상수인 특수한 경우)도 유리함수에 포함됨을 확인 $\to$ 문제 1.\newline
\textbf{정의역 구하기 (8분)} --- 분모 $\neq 0$인 $x$ 전체가 정의역임을 확인 $\to$ 문제 2.\newline
\textbf{식의 변형 (12분)} --- $y=\dfrac{ax+b}{cx+d}$를 다항식의 나눗셈으로 $y=\dfrac{k}{x-p}+q$ 꼴로 변형하는 절차(분자를 분모로 나눈 몫과 나머지 이용) $\to$ 예제, 문제 3, 문제 4. &
ATT: 촉진자 역할\newline ATL: 대수적 조작 \\
\addlinespace
\textcolor{goldstrong}{\textbf{정리}}\newline\footnotesize 10분 &
\textbf{개념 연결 질문 (5분)} --- ``$y=\dfrac{k}{x-p}+q$ 꼴로 바꾸면 무엇이 편리해질까?'' (다음 시간 그래프 예고)\newline
\textbf{성찰 (5분)} --- 감사 일기 1문장 + 다음 차시 예고(유리함수의 그래프). &
마음공부 · 감사 일기 \\
\bottomrule
\end{tabularx}

\begin{calloutbox}
\textbf{지도상의 유의점} --- $y=\dfrac{ax+b}{cx+d}$ 꼴을 $y=\dfrac{k}{x-p}+q$ 꼴로 바꾸는 과정에서 분자를 분모로 나누는 다항식의 나눗셈을 어려워하는 학생이 많다. $ax+b = a\cdot\frac{c}{c}\left(x+\frac{d}{c}\right) - \frac{ad}{c}+b$ 형태로 분모 $cx+d$가 드러나도록 분자를 인위적으로 조작하는 과정을 단계별로 짚어 준다.
\end{calloutbox}

\section{형성평가 \& 답안}
\begin{minipage}[t]{0.48\linewidth}
\begin{answerbox}
\textbf{문제 1} \quad 다음 중 유리함수인 것을 모두 고르시오.\\
\small $y=2x+1$, \; $y=\dfrac1x$, \; $y=x^2-3$, \; $y=\dfrac{x+1}{x-2}$, \; $y=\sqrt x$\\[4pt]
\textbf{\color{goldstrong}답} \; $y=\sqrt x$를 제외한 모두. (다항함수도 분모가 $1$인 유리함수)
\end{answerbox}
\end{minipage}\hfill
\begin{minipage}[t]{0.48\linewidth}
\begin{answerbox}
\textbf{문제 2} \quad 유리함수 $y=\dfrac{3}{x-2}$의 정의역을 구하시오.\\[4pt]
\textbf{\color{goldstrong}답} \; $\{x\mid x\neq 2$인 실수$\}$
\end{answerbox}
\end{minipage}

\vspace{6pt}
\begin{answerbox}
\textbf{문제 3} \quad $y=\dfrac{2x+1}{x-1}$을 $y=\dfrac{k}{x-p}+q$ 꼴로 변형하고 정의역을 구하시오.\\[4pt]
\textbf{\color{goldstrong}답} \; $2x+1=2(x-1)+3$이므로 $y=\dfrac{3}{x-1}+2$. 정의역: $\{x\mid x\neq1$인 실수$\}$
\end{answerbox}
\begin{answerbox}
\textbf{문제 4} \quad $y=\dfrac{3x-5}{x+2}$를 $y=\dfrac{k}{x-p}+q$ 꼴로 변형하고 정의역을 구하시오.\\[4pt]
\textbf{\color{goldstrong}답} \; $3x-5=3(x+2)-11$이므로 $y=-\dfrac{11}{x+2}+3$. 정의역: $\{x\mid x\neq-2$인 실수$\}$
\end{answerbox}

\section{성취수준 (이 차시 범위)}
\begin{tabularx}{\linewidth}{@{}>{\bfseries\color{goldstrong}}p{0.9cm}X@{}}
\toprule
A & 유리함수와 다항함수의 관계를 설명하고, 복잡한 유리식도 $\dfrac{k}{x-p}+q$ 꼴로 정확히 변형하며 정의역을 구할 수 있다. \\
B & 유리함수의 정의역을 구하고, 식을 $\dfrac{k}{x-p}+q$ 꼴로 변형할 수 있다. \\
C & 안내에 따라 간단한 유리함수의 정의역을 구할 수 있다. \\
D & 유리함수와 다항함수를 구별할 수 있다. \\
E & 유리식이 무엇인지 안다. \\
\bottomrule
\end{tabularx}

\section{다음 차시}
\begin{calloutbox}
\textbf{01 유리함수 --- 2/2차시.} $y=\dfrac{k}{x}$의 그래프와 평행이동으로 만들어지는 유리함수의 그래프, 점근선을 다룬다.
\end{calloutbox}
\end{document}
